\section{The nonzero integral kernel class}
\label{sec:main-proof}

The proof now has all three inputs: the integral cyclotomic class,
the divisor $\Spec\OO\subset\Spec A$ and its four split points, and
the specialization criterion. We first push forward integrally to
obtain a class killed by generic restriction. We then detect its
coefficient reduction after flat pullback by the elliptic orbit.

\begin{proof}[Proof of Theorem~\ref{thm:main}]
Let $c\in K_5(S)$ be the class from Proposition~\ref{prop:integral-lift},
and write $c_L\in K_5(L)$ for its restriction along $S\to L$.
There is no map $S\to\OO$ in this construction, since $5$ is invertible
in $S$. Instead, integral localization for
the DVR $\OO$ gives
\[
 K_5(\OO)\longrightarrow K_5(L)\longrightarrow K_4(k)=0,
\]
where the last equality is Quillen's computation for finite fields
\cite{Quillen1972FiniteFields}; see also \cite[IV.1.13]{Weibel2013}.
The residue field here has characteristic five; the equality is the
integral finite-field calculation. Choose an integral lift
$\widetilde c\in K_5(\OO)$. The closed immersion
$i:\Spec\OO\hookrightarrow\Spec A$ of Lemma~\ref{lem:local-ring} defines
\[
 \gamma=i_*\widetilde c\in G_5(A)=K_5(A).
\]
The equality uses the canonical resolution isomorphism for the regular
ring $A$. Restriction to $\operatorname{Frac}A$ annihilates every coherent
module supported on $\Spec\OO$. Thus $\gamma$ is in the kernel of the
canonical augmentation, integrally.

To show that it is nonzero, make the flat base change $V\to L$. The square
\[
\begin{tikzcd}[column sep=large,row sep=large]
 \displaystyle\coprod_{\delta\in\Delta}\Spec L
       \arrow[r,"\coprod p_\delta"] \arrow[d]
       & T \arrow[d] \\
 \Spec\OO \arrow[r,"i"] & \Spec A
\end{tikzcd}
\]
is Cartesian by Lemma~\ref{lem:splitting}. To check flat base change with the point indices, put $B_T=A\otimes_VL$.
For every finitely generated $\OO$-module $M$, scalar extension gives
natural isomorphisms of $B_T$-modules
\begin{equation}\label{eq:module-base-change}
 \begin{aligned}
 B_T\otimes_A i_*M
 &\simeq(\OO\otimes_VL)\otimes_{\OO}M\\
 &\simeq\bigoplus_{\delta\in\Delta}L\otimes_{\OO,\delta}M.
 \end{aligned}
\end{equation}
The $\delta$-summand is a $B_T$-module through the quotient defining
$p_\delta$. These are exact functors: $B_T$ is flat over $A$, and
each map $\OO\to L$ is a localization followed by an automorphism.
Thus \eqref{eq:module-base-change} proves the required identity on
integral coherent $K$-theory.

Write $\iota:S\to L$ for the fixed embedding. Since $\widetilde c$
restricts to $c_L=\operatorname{res}_\iota c$, its image on the
$\delta$-factor is
\[
 \operatorname{res}_{\delta:\OO\to L}\widetilde c
 =\operatorname{res}_{\iota\circ\delta}c
 =\operatorname{res}_{\iota}(\delta c).
\]
The equality uses the specified identification
$\Delta=\Gal(L/H_0)$ and requires no equivariance of the chosen DVR
preimage. Consequently, still integrally,
\[
 \gamma|_T=\sum_{\delta\in\Delta}(p_\delta)_*
 \operatorname{res}_\iota(\delta c).
\]
Naturality of coefficient reduction and
Proposition~\ref{prop:integral-lift} now give
\begin{equation}\label{eq:gamma-split}
 \overline{\gamma|_T}
 =\sum_{\delta\in\Delta}(p_\delta)_*
       \bigl(\chi(\delta)^2b^2[\delta u]_5\bigr)
       \quad\text{in }G_5(T;5).
\end{equation}
Here $b$ uses the fixed embedding $S\to L$, and the character factors
come from Proposition~\ref{prop:integral-lift}. No choice of the integral
lift $\widetilde c$ affects this formula.

Apply Proposition~\ref{prop:specialization} to the model $\Dcal/\OO$,
to $U=T$, and to the set $Z_\infty$ and homomorphism $\psi$ constructed
in Section~\ref{sec:finite-geometry}. The required
specialization condition follows directly from the definition of $A$:
let $q$ be any closed point of $C$ and let $\OO_q$ be the valuation
ring of its finite residue-field extension $\kappa(q)/L$. If $q$
specializes in the chart $Z\ne0$, the proper extension
$\Spec\OO_q\to\Dcal$ lies in that chart, since an open set containing
the closed point of a DVR contains its generic point. Hence $X/Z$ and
$Y/Z$ belong to $\OO_q$. The original coordinates
$x=tX/Z$ and $y=tY/Z$ have positive valuation. Every denominator inverted
outside $(5,x,y)$ has nonzero constant residue and hence is a unit at $q$.
Thus $q$ belongs to $T$.

All the hypotheses of Proposition~\ref{prop:specialization} are now
verified: the smooth model and localized chart come from
Section~\ref{sec:local-geometry}, the retained-point condition holds
for every closed point, and \eqref{eq:detector-at-infinity} says that
$\psi$ kills the entire set $Z_\infty$.
The points $p_\delta$ specialize to $g_{\chi(\delta)}P$ and are
$L$-rational. Every $\delta u$ is a uniformizer, so
$v_L(\delta u)=1$. The sum tested by
Proposition~\ref{prop:specialization} for~\eqref{eq:gamma-split} is therefore
\[
 \sum_{\delta\in\Delta}\chi(\delta)^2
        \psi([g_{\chi(\delta)}P])
 =4Q\ne0\quad\text{in }E(k)/5E(k),
\]
by Lemma~\ref{lem:orbit}. It follows that
$\overline{\gamma|_T}\ne0$. Hence $\gamma\ne0$ in the original integral
group, as required.
\end{proof}

\begin{remark}
Although the divisor $i$ is principal, its pushforward need not vanish on
all higher $K$-groups. The self-intersection formula gives $i^*i_*=0$ here;
the projection formula gives $i_*i^*=0$. Neither formula asserts $i_*=0$.
In particular, the proof does not assume that $\widetilde c$ is the
restriction of a class on $A$.
\end{remark}


The positions in the signed orbit matter for this detector. Forgetting
them and retaining only the four normalized valuations, all equal to one,
gives the scalar sum $1+4+4+1=0$ in $\F_5$. The Picard test retains the
points and instead gives $4Q$. This calculation describes that particular
scalar forgetting; it makes no assertion about all scalar functionals.
